\par\addvspace{12pt}\Needspace{9\baselineskip}

\section{极点极线}

% ITEM: quadratic-example-079
\solutionitem{quadratic-example-079}{例4.79}
对圆\(x^2+y^2=4\)，分别求\((1,0)\)、\((2,0)\)、\((4,0)\)的极线，比较它们与圆的位置关系。再求直线\(x+y=1\)关于椭圆\(x^2/9+y^2/4=1\)的极点。
\solutionanswer
三点极线依次为\(x=4,x=2,x=1\)，分别与圆相离、相切、相交。椭圆的极点由\(ux/9+vy/4=1\)匹配\(x+y=1\)，得到\((9,4)\)。


% ITEM: quadratic-example-080
\solutionitem{quadratic-example-080}{例4.80}
求直线\(2x-3y+4=0\)关于抛物线\(y^2=4x\)的极点。设该直线与抛物线交于\(A,B\)，求两点处切线的交点，并核对其与极点的关系。
\solutionanswer
抛物线极线\(vy=2(x+u)\)与给定直线匹配，极点为\((2,3)\)。直线与曲线交于\((1,2),(4,4)\)，两切线\(y=x+1,y=x/2+2\)交于\((2,3)\)，与极点一致。


% ITEM: quadratic-exercise-46-01
\solutionitem{quadratic-exercise-46-01}{习题4.46（1）}
分别求\((1,1)\)关于圆\(x^2+y^2=1\)、椭圆\(x^2/4+y^2=1\)和双曲线\(x^2-y^2=1\)的极线。
\par
\solutionanswer
由标准曲线极线公式，对圆\(x^2+y^2=1\)，点\((1,1)\)的极线为\(x+y=1\)；对椭圆\(x^2/4+y^2=1\)，为\(x/4+y=1\)；对双曲线\(x^2-y^2=1\)，为\(x-y=1\)。三式均为普通平面中的直线。


% ITEM: quadratic-exercise-46-02
\solutionitem{quadratic-exercise-46-02}{习题4.46（2）}
求\(x+2y=3\)关于椭圆\(x^2/9+y^2/4=1\)的极点，并判断该极点能否向椭圆作两条实切线。
\par
\solutionanswer
将直线化为\(x/3+2y/3=1\)，与椭圆极线公式\(ux/9+vy/4=1\)比较，得到极点\(P=(3,8/3)\)。其椭圆二次式值为
\[\frac{3^2}{9}+\frac{(8/3)^2}{4}=\frac{25}{9}>1,\]
所以点在椭圆外，能向椭圆作两条不同的实切线；该点的极线正是两切点连线所在直线。


% ITEM: quadratic-exercise-46-03
\solutionitem{quadratic-exercise-46-03}{习题4.46（3）}
求点\((u,v)\)关于圆\((x-h)^2+(y-k)^2=R^2\)的极线。由此研究圆心与过圆心的直线在极点极线对应中的情形。
\par
\solutionanswer
移到圆心坐标，极线为\((u-h)(x-h)+(v-k)(y-k)=R^2\)。极点为圆心时左边恒为零，不表示普通直线。任何有限点的极线都不经过圆心，因此过圆心直线没有有限极点。若直线为\(\alpha(x-h)+\beta(y-k)=c\ne0\)，则唯一极点满足\((u-h,v-k)=R^2(\alpha,\beta)/c\)。


% ITEM: quadratic-exercise-46-04
\solutionitem{quadratic-exercise-46-04}{习题4.46（4）}
对抛物线\(y^2=8x\)，分别求直线\(x-y+1=0\)、\(x=2\)的极点；说明\(y=2\)为什么没有有限极点。
\par
\solutionanswer
抛物线\(y^2=8x\)对应\(p=4\)。其极线公式为\(vy=4(x+u)\)；直线\(\alpha x+\beta y+\gamma=0\)（\(\alpha\ne0\)）的极点为\((\gamma/\alpha,-4\beta/\alpha)\)。
对\(x-y+1=0\)，极点为\((1,4)\)，其极线\(4y=4(x+1)\)正好等价于原直线。对\(x=2\)，极点为\((-2,0)\)，其极线\(0=4(x-2)\)即\(x=2\)。直线\(y=2\)的\(x\)系数为零，不能按有限点极线中的方程\(p x-vy+pu=0\)匹配出有限\((u,v)\)：该式的\(x\)系数恒为\(p=4\ne0\)，故它没有有限极点。


% ITEM: quadratic-exercise-46-05
\solutionitem{quadratic-exercise-46-05}{习题4.46（5）}
设\(P\)在圆\(x^2+y^2=R^2\)外，极线与\(OP\)交于\(H\)。证明\(|OP|\,|OH|=R^2\)，并比较\(H\)与两切点的弦中点。
\par
\solutionanswer
令\(P\)的位置向量为\(p\)，则其极线为\(p\cdot x=R^2\)。交点\(H\)在\(OP\)上，写成\(h=\lambda p\)，代入极线得\(\lambda|p|^2=R^2\)，所以\(|OH|=\lambda|OP|=R^2/|OP|\)，即\(|OP|\,|OH|=R^2\)。
若\(P\)在圆外，两条切线切点为\(T_1,T_2\)。两点均满足\(P\cdot T_i=R^2\)，故弦\(T_1T_2\)所在直线就是极线。由于\(OT_i=R\)，弦中点到圆心的连线垂直于弦；而\(OP\)也垂直于极线，所以弦中点是\(OP\)与极线的交点\(H\)。


% ITEM: quadratic-exercise-46-06
\solutionitem{quadratic-exercise-46-06}{习题4.46（6）}
对\(xy=1\)，由定义求点\((u,v)\)的极线。分别研究\(uv=0\)与\(uv\ne0\)时极线和曲线的实交点，并与上一节的过点切线条数比较。
\par
\solutionanswer
对\(\Phi(x,y)=xy-1\)，极化后点\((u,v)\)的极线为
\[
\frac{uy+vx}{2}-1=0,\qquad\text{即 }vx+uy=2.
\]
若\(uv\ne0\)，以\(y=(2-vx)/u\)代入\(xy=1\)，得到\(v x^2-2x+u=0\)，判别式为\(4(1-uv)\)。因此\(uv<1\)时有两个不同实交点，\(uv=1\)时有一个切点（重根），\(uv>1\)时无实交点。
若\(u=0,v\ne0\)，极线为\(x=2/v\)，与双曲线有唯一交点\((2/v,v/2)\)；若\(v=0,u\ne0\)，极线为\(y=2/u\)，唯一交点为\((u/2,2/u)\)。若\(u=v=0\)，极化式成为\(-1=0\)，并不表示一条有限极线，原点是双曲线中心。

过\((u,v)\)的切线在接触点\((t,1/t)\)处满足\(x/t+ty=2\)。令其经过\((u,v)\)，得\(vt^2-2t+u=0\)，与上述交点判别具有相同的判别式；不过参数必须满足\(t\ne0\)。当\(uv\ne0\)时，切线条数分别为两条、一条、零条，对应\(uv<1,uv=1,uv>1\)。当\(uv=0\)且\((u,v)\ne(0,0)\)时，二次式降次，或形式根\(t=0\)不属于双曲线参数，故只有一条实切线，与极线和双曲线唯一交点相应。中心\((0,0)\)没有实切线，且没有有限极线。


% ITEM: quadratic-exercise-46-07
\solutionitem{quadratic-exercise-46-07}{习题4.46（7）}
验证式\textup{(4.155)}。再将\(x^2/4+y^2=1\)作变换\(X=x+y+1,Y=y-1\)，求点\((2,1)\)及其极线的像，核对像曲线的极线方程。
\par

\par\smallskip\noindent 正文所引关系为：
\begin{equation*}
T(P,X)=2\Phi\left(\frac{P+X}{2}\right)
       -\frac{\Phi(P)+\Phi(X)}{2}.
\end{equation*}
\solutionanswer
将\(\Phi=q+L^TX+F\)代入右边并展开，平方交叉项留下\(b(P,X)\)，一次项留下\(L^T(P+X)/2\)，常数项留下\(F\)，正是\(T(P,X)\)。给定点映为\((4,0)\)，原极线\(x/2+y=1\)映为\(X+Y=2\)。像曲线为\((X-Y-2)^2/4+(Y+1)^2=1\)，在其中按相同极化公式求\((4,0)\)的极线，也得到\(X+Y=2\)。


% ITEM: quadratic-example-081
\solutionitem{quadratic-example-081}{例4.81 *}
对\(x^2+2xy+2y^2-2x=0\)，写出完整矩阵，求\((1,1)\)的极线与直线\(y=1\)的极点。再找出没有普通极线的有限点，并求所有没有有限极点的直线。
\solutionanswer
完整矩阵\(M=\begin{bmatrix}1&1&-1\\1&2&0\\-1&0&0\end{bmatrix}\)。\((1,1)\)的极线为\(x+3y=1\)，\(y=1\)的极点为\((1,0)\)。有限点\((2,-1)\)的极线式为非零常数，故没有普通极线。对直线\(\alpha x+\beta y+\gamma=0\)，\(M^{-1}l\)末坐标为\((\beta-\gamma-2\alpha)/2\)，因此没有有限极点的直线恰为所有经过\((2,-1)\)的直线。


% ITEM: quadratic-example-082
\solutionitem{quadratic-example-082}{例4.82}
点\(P\)在直线\(x+2y=3\)上移动，求它关于椭圆\(x^2/9+y^2/4=1\)的极线经过的定点。若改为\(x+2y=0\)，求这些极线的公共方向。
\solutionanswer
给定直线的极点为\((3,8/3)\)，由互反性，所有动极线都经过它。若动点线改为\(x+2y=0\)，非零动点的极线为\(-8x+9y=\text{常数}\)，公共方向为\((9,8)\)；中心本身不产生普通极线。


% ITEM: quadratic-example-083
\solutionitem{quadratic-example-083}{例4.83}
对抛物线\(y^2=4x\)，点\(Q=(1,0)\)。求轴上点\(P\)，使过\(Q\)的各条割线所截端点\(M,N\)满足\(k_{PM}+k_{PN}=0\)。用移轴齐次化求解，再验证\(Q\)在\(P\)的极线上，并解释相应的等角关系。
\solutionanswer
设\(P=(p,0)\)。若\(p=1\)，观察点就是\(Q\)，一般斜率为\(k\)的割线会给出斜率和\(2k\)，不合要求。以下设\(p\ne1\)。
移轴\(u=x-p,v=y\)，抛物线为\(v^2-4u-4p=0\)。过\(Q\)的非竖直割线写成\(v=k(u+p-1)\)，其中\(k\ne0\)。利用\(v-ku=k(p-1)\)齐次化，得到
\[
 \bigl[k^2(p-1)^2-4p\bigr]v^2
 +4k(p+1)uv-4k^2u^2=0.
\]
故两个观察斜率\(r=v/u\)满足同一二次方程，其和为
\[
 -\frac{4k(p+1)}{k^2(p-1)^2-4p}.
\]
要对所有有意义的割线恒为零，必有\(p=-1\)。代回后两根互为相反数；竖直割线\(x=1\)给出斜率\(\pm1\)，也满足要求。
因此\(P=(-1,0)\)，其极线为\(x=1\)，确经过\(Q\)。两观察斜率相反，表示两连线关于对称轴等角。

另可先从割线与抛物线的联立方程得到\(y_1y_2=-4\)。斜率和的分子为\(-(p+1)(y_1+y_2)\)，也能直接得到\(p=-1\)。移轴齐次化则把同一条件直接写成两条观察方向的关系。


% ITEM: quadratic-example-084
\solutionitem{quadratic-example-084}{例4.84}
对椭圆\(x^2/9+y^2/4=1\)，求过右焦点且斜率为\(1\)的弦两端切线的交点。再改变弦的斜率，说明这些交点所在的直线，并判断焦点所在轴的情形。
\solutionanswer
弦线为\(x-y=\sqrt5\)，匹配极线公式得交点\((9/\sqrt5,-4/\sqrt5)\)。随焦弦方向改变，有限交点都在右准线\(x=9/\sqrt5\)上。焦点轴上的弦连接两顶点，其切线平行，没有有限交点，须与一般情形区分。


% ITEM: quadratic-example-085
\solutionitem{quadratic-example-085}{例4.85}
回到例4.75，将点\(D\)所在直线看成抛物线的准线，仅用焦点与准线的极点极线关系证明切点弦经过定点。再用坐标方法核对定点。
\solutionanswer
\(x^2=2y\)的焦点\(F=(0,1/2)\)，其极线恰为准线\(y=-1/2\)。\(D\)位于该极线上，互反性使\(D\)的极线经过\(F\)，而\(D\)的极线就是切点弦\(AB\)。坐标上令\(D=(d,-1/2)\)，得到\(y=dx+1/2\)，固定点同为\((0,1/2)\)。


% ITEM: quadratic-exercise-47-01
\solutionitem{quadratic-exercise-47-01}{习题4.47（1）}
点\(P=(2,1)\)关于圆\(x^2+y^2=1\)的极线为\(\ell\)。任取\(Q\in\ell\)，写出其极线，并直接验证它经过\(P\)。
\par
\solutionanswer
若\(Q=(u,v)\)，则\(2u+v=1\)。其极线\(ux+vy=1\)代入\(P=(2,1)\)即成立，故经过\(P\)。


% ITEM: quadratic-exercise-47-02
\solutionitem{quadratic-exercise-47-02}{习题4.47（2）}
设\(A,B,C\)是椭圆上三个不同点，三条端点切线两两交于\(P,Q,R\)。证明三条边\(AB,BC,CA\)分别是相应切线交点的极线。若\(P,Q,R\)均存在，它们能否共线？
\par
\solutionanswer
两切线交点位于两个切点的切线上，由互反性，两切点都在交点的极线上，所以该极线正是相应弦。三个交点不能共线：三条不同切线的两两交点若共线，三切线只能共点，而从一点向椭圆至多能作两条实切线，矛盾。


% ITEM: quadratic-exercise-47-03
\solutionitem{quadratic-exercise-47-03}{习题4.47（3）}
求椭圆\(x^2/16+y^2/9=1\)中斜率为\(1\)的平行弦中点所在直径，以及各弦端点切线交点所在直线。若某弦中点为\((4/5,-9/20)\)，求切线交点。
\par
\solutionanswer
中点直径为\(9x+16y=0\)。各平行弦的极点也沿这条直线，因为其极线法向量比例固定；中心弦的两切线平行，没有有限交点。对给定\(M\)，\(q(M)=1/16\)，故切线交点\(S=M/q(M)=(64/5,-36/5)\)。由\(M^TH(X-M)=0\)与\(S^THX=1\)也可直接核对。


% ITEM: quadratic-exercise-47-04
\solutionitem{quadratic-exercise-47-04}{习题4.47（4）}
设抛物线\(y^2=2px\)上一条弦的端点切线交于\(S\)，弦中点为\(M\)。证明\(S_y=M_y\)，并解释平行弦的中点轨迹及端点切线交点轨迹之间的关系。
\par
\solutionanswer
切点参数\(t_1,t_2\)给出\(S=(pt_1t_2/2,p(t_1+t_2)/2)\)，中点的纵坐标也是\(p(t_1+t_2)/2\)。所以\(S_y=M_y\)。平行弦的参数和固定，两类轨迹都在同一条水平直线上；横坐标还满足\(S_x+M_x=M_y^2/p\)，区分实际取值范围。


% ITEM: quadratic-exercise-47-05
\solutionitem{quadratic-exercise-47-05}{习题4.47（5）}
椭圆\(x^2/25+y^2/9=1\)的一条焦弦两端切线交于\(S\)，且\(S\)在直线\(y=2x\)上。求\(S\)及焦弦方程。
\par
\solutionanswer
两条准线为\(x=\pm25/4\)，焦弦两端切线交点在对应准线上。联立\(y=2x\)得\(S=(25/4,25/2)\)或其相反点。两条焦弦是相应极线，方程为\(x/4+25y/18=\pm1\)。


% ITEM: quadratic-exercise-47-06
\solutionitem{quadratic-exercise-47-06}{习题4.47（6）}
抛物线\(y^2=4x\)的焦弦\(AB\)两端切线交于\(S\)。证明\(S\)在准线上且\(SA\perp SB\)。若\(|SA|=|SB|\)，求\(S\)和\(AB\)。
\par
\solutionanswer
焦点的极线是准线\(x=-1\)，故\(S\)在其上。切点参数乘积为\(-1\)，两切线斜率之积亦为\(-1\)，所以垂直。若切线段等长，由\(|SA|=|t_1-t_2|\sqrt{1+t_1^2}\)及对\(B\)的对应式，得\(t_1^2=t_2^2\)，进而\(t_1=-t_2=\pm1\)。因此\(S=(-1,0),AB:x=1\)。


% ITEM: quadratic-exercise-47-07
\solutionitem{quadratic-exercise-47-07}{习题4.47（7）}
设\(P\)在椭圆\(x^2/a^2+y^2/b^2=1\)外，切点弦为\(AB\)。取\(Q\)在线段\(AB\)内部，从\(Q\)还能否向椭圆作实切线？\(Q\)的极线与\(P\)有什么关系？说明这些结论如何同时成立。
\par
\solutionanswer
\(Q\)在两切点之间的开弦内，严格位于椭圆内部，不能作实切线；但仍有极线。因\(Q\)在\(P\)的极线上，互反性保证\(Q\)的极线经过\(P\)。这条极线与椭圆相离，经过一个外点并不要求它与椭圆相交，两结论并无冲突。


% ITEM: quadratic-exercise-47-08
\solutionitem{quadratic-exercise-47-08}{习题4.47（8）}
对双曲线\(x^2/9-y^2/4=1\)，求右焦点、右准线对应的极线和极点。取过右焦点的一条实焦弦，若两端切线相交，证明其交点在右准线上，并说明何种方向只产生一个有限交点。
\par
\solutionanswer
右焦点\((\sqrt{13},0)\)的极线为右准线\(x=9/\sqrt{13}\)，该准线的极点反过来就是右焦点。焦弦端点切线的交点是焦弦的极点，因而在准线上。平行于\(y=\pm2x/3\)的直线只产生一个有限交点；焦点轴上的弦虽有两个端点，但两端切线平行，不存在有限交点。


% ITEM: quadratic-exercise-47-09
\solutionitem{quadratic-exercise-47-09}{习题4.47（9）}
过椭圆内一个非中心点\(P\)作两条弦\(AB,CD\)，令\(U=AC\cap BD,V=AD\cap BC\)，设交点均存在。证明\(U,V\)在\(P\)的极线上。可先圆化，再由坐标证明。若\(P\)取为焦点，指出\(U,V\)所在直线，并说明如何将这个结论用于抛物线的两条焦弦。
\par
\solutionanswer
以\(P\)为原点、两弦为两坐标轴，设端点\(A=(\alpha,0),B=(\beta,0),C=(0,\gamma),D=(0,\delta)\)，均不为原点。方程记为\(ax^2+2hxy+by^2+2gx+2fy+c=0\)。于是\(a\alpha\beta=b\gamma\delta=c\)，\(a(\alpha+\beta)=-2g,b(\gamma+\delta)=-2f\)。\(U=AC\cap BD\)满足两截距式，相加得
\[
(1/\alpha+1/\beta)x_U+(1/\gamma+1/\delta)y_U=2,
\]
即\(gx_U+fy_U+c=0\)，正是\(P\)的极线。交换\(C,D\)得到\(V\)的同一结论。\(P\)为焦点时该极线是相应准线；同样证明也适用于抛物线的两条焦弦，两个交叉连接点均在准线上。
\par
